\section{方法：残差学习}

残差学习模块（Residual Block）是ResNet的核心构建单元。设模块的输入为 $x$，期望的基础映射为 $H(x)$，则残差映射定义为 $F(x) = H(x) - x$。那么原始映射可表达为 $H(x) = F(x) + x$。其中 $F(x)$ 由一系列堆叠的卷积层、批量归一化层和激活函数层实现，而“$+x$”通过一个恒等快捷连接（Identity Shortcut Connection）完成逐元素加法。这种结构具有以下关键设计原则：

\textbf{快捷连接不增加额外参数或计算量}。恒等快捷连接仅需逐元素加法，其梯度可以直接绕过中间的权重层反向传播，从而缓解梯度消失。同时，残差框架不限制网络的深度，因为 $F(x)$ 的复杂度可以任意设计。

\textbf{为什么残差更易优化？} 从优化观点看，当最优映射接近恒等时，残差 $F(x)$ 趋向于一个小的扰动，即 $F(x) \approx 0$。此时网络只需学习如何使输出趋近于零，而堆叠的非线性层拟合零函数远容易于拟合恒等函数。此外，在随机初始化时权重通常接近于零，因此 $F(x)$ 的初始输出也接近零，使得 $H(x) \approx x$，这为网络提供了一个良好的初始状态。

\textbf{维度匹配与调整}。当输入 $x$ 的维度与 $F(x)$ 的维度不同时（例如跨层下采样导致特征图尺寸减半或通道数改变），需要调整 $x$ 的维度以进行逐元素加法。常用的方法有三种：\textbf{补零（Zero-padding）}，在特征图边缘填充零以对齐尺寸；\textbf{投影（Projection）}，使用额外的 $1\times1$ 卷积层将 $x$ 投射到目标维度（通常步长为2以同时下采样）；\textbf{$1\times1$ 卷积}，专门用于通道数匹配。实验表明，恒等映射（无额外参数）的效果优于投影，因此在推荐设计中使用恒等映射，仅在维度变化处使用 $1\times1$ 卷积。

残差块的具体实现有两种常用形式：
\begin{itemize}
    \item \textbf{基础块（Basic Block）}：用于较浅的网络（ResNet-18/34），由两个 $3\times3$ 卷积层堆叠，每个卷积后接批量归一化和ReLU激活，最后与输入相加。
    \item \textbf{瓶颈块（Bottleneck Block）}：用于深层网络（ResNet-50/101/152），通过 $1\times1 \rightarrow 3\times3 \rightarrow 1\times1$ 的结构降低计算量。第一个 $1\times1$ 降维（例如 256→64），$3\times3$ 卷积在低维上运算，最后一个 $1\times1$ 升维回256。瓶颈块在相同参数量下可构建更深的网络。
\end{itemize}

通过堆叠这些残差块并配合阶段性下采样，ResNet可以轻松扩展到超过100层的深度，而训练难度与传统30层网络相当。